\chapter{Discussion and Conclusion (WIP)}
\label{ch:5}
\label{sec:5}


This chapter will interpret the results reported in \cref{sec:4} without introducing new evidence. It will address the ceiling of seven-class EPC prediction, the unit-of-observation problem in EPC labels, the appropriate role of SVI in multimodal building-data enrichment, the study limitations, future work, and the final bounded conclusion.


\section{Observations (WIP)}
\label{sec:5_1}


\subsection{Seven-Class Ceiling: Task Structure Rather Than Visual Failure (WIP)}
\label{sec:5_1_1}


\begin{quote}
WIP summary. The proposed interpretation attributes the low seven-class ceiling to four interacting constraints: limited information in the available attributes, biased pool composition, unobserved building-system information, and inconsistent building-level EPC labels. The small gap between the reference-attribute and attribute-mediated routes suggests that visual substitution is not the sole cause of limited fine-grained performance.
\end{quote}


\subsection{Unit-of-Observation Mismatch in EPC Labels (WIP)}
\label{sec:5_1_2}


\begin{quote}
WIP summary. EPCs may be issued at dwelling-unit level while the model predicts at building level. The discussion will interpret within-building certificate disagreement as a failure of the assumed building-level homogeneity and recommend either unit-level prediction or an explicitly defined and sensitivity-tested building-level aggregation rule.
\end{quote}


\subsection{Boundary of the Role of SVI: Division of Labour Between Modalities (WIP)}
\label{sec:5_1_3}


\begin{quote}
WIP summary. The working principle is that SVI should supply semantic attributes, whereas reconstructed 3D open data should supply geometry. Attribute-mediated VLM inference appears more effective than direct EPC prompting for the evaluated InternVL3-2B configuration, but this result is bounded to the tested model and does not establish a universal rule for all VLMs.
\end{quote}


\section{Limitations (WIP)}
\label{sec:5_2}


\begin{quote}
WIP summary. The principal limitations are contributor-driven image-coverage bias, imperfect building-level EPC labels, evaluation of cross-city transfer using only Amsterdam as the target city, validation in a single country, operational simplifications in the TABULA mapping, incomplete coverage of model paradigms, and use of \(H'_{tr}\) as an internal physical reweighting rather than validation against measured energy consumption.
\end{quote}


\section{Future Work (WIP)}
\label{sec:5_3}


\begin{quote}
WIP summary. Priorities include testing image-based renovation cues through residual proxy labels, repeating leave-one-city-out evaluation for the remaining cities, collecting targeted suburban imagery, developing cross-national archetype mappings, and moving towards apartment-floor or dwelling-unit prediction.
\end{quote}


\section{Conclusion}
\label{sec:5_4}


This study evaluated whether SVI-derived building attributes can substitute for reference attributes in TABULA-based building-stock enrichment. A value-traceable semi-synthetic dataset linked attribute extraction, archetype assignment, physical parameters, and downstream EPC classification within a common comparison framework. The results provide clear but scope-bounded answers to the three research objectives.


For Objective 1, SVI supported the extraction of building type, construction year, and floor count and therefore served as an observational source for missing semantic building information. The frozen DINOv2 backbone provided the most balanced performance across the three attributes, ResNet-50 had a slight advantage only in building-type accuracy, and zero-shot InternVL3-2B was substantially weaker. These findings show that frozen self-supervised visual representations can provide effective features for multi-attribute building-profile enrichment.


For Objective 2, the extracted attributes had task-specific rather than uniform downstream utility. Construction year contributed the most information to EPC classification, while floor count made a smaller but directionally consistent contribution. Building type was structurally necessary for TABULA cell assignment but provided limited additional EPC discrimination in the apartment-dominated sample. Joint-cell evaluation further showed that strong individual-attribute performance primarily translated into reliable assignment of dominant archetype cells rather than balanced performance across the full TABULA matrix. The analysis therefore identified not only which attributes could be extracted, but also which attributes should be prioritised and which downstream tasks benefited from them.


For Objective 3, SVI-derived attributes provided an informative partial substitute when reference metadata were missing, but they were not equivalent to reference attributes. The strongest attribute-mediated route retained discrimination above the no-information baseline in binary EPC classification, with a ROC-AUC substitution cost of 0.065 relative to the reference-attribute route. The physical-consequence analysis likewise showed that the trained configurations controlled archetype-derived heat-loss error more effectively than the zero-shot configuration. Cross-city evaluation, however, showed that substitution performance depended on alignment between the training and deployment distributions and may require target-domain calibration or more representative training coverage. This transfer boundary defines the conditions under which the approach applies without negating the information supplied by SVI under distribution-matched conditions.


Overall, SVI-derived attributes constitute a conditional enrichment source for buildings with missing reference metadata. Construction year provided the clearest downstream utility, while the frozen self-supervised representation offered the best overall balance between upstream extraction and control of physical consequences. The main contribution is a traceable evidence chain that separately quantifies which attributes are useful, the cost of substituting their source, and the conditions under which the substitution remains applicable. Within the Netherlands sample, which is dominated by Amsterdam apartment blocks, SVI-derived attributes retained a meaningful share of downstream information and supported interpretable TABULA-based building-stock enrichment.
